\documentclass[12pt,a4paper]{article}

% 這是一份latex 文件的 preamble
% 請AI閱讀後,將附上的PDF整理,以latex code整理,儘量用longtable
% 請注意我的名字: 謝慕揚, MD, PhD, FESC
% 表格請注意換行,不該在cell內使用 \\
% 表格請不要有垂直分隔線 |

% Including essential packages for Chinese typesetting
\usepackage{xeCJK}
\usepackage{fontspec}
% Configuring Chinese font with Noto Serif CJK TC, fallback to SimSun
\IfFontExistsTF{Noto Serif CJK TC}{
  \setCJKmainfont{Noto Serif CJK TC}
  \setCJKsansfont{Noto Serif CJK TC}
  \setCJKmonofont{Noto Serif CJK TC}
}{\setCJKmainfont{SimSun}
  \setCJKsansfont{SimSun}
  \setCJKmonofont{SimSun}
}

% Including other necessary packages
\usepackage{geometry}
\usepackage{titlesec}
\usepackage{enumitem}
\usepackage{xcolor}
\usepackage{colortbl}
\usepackage{tcolorbox}
\usepackage{booktabs}
\usepackage{longtable}
\usepackage{array}
\usepackage{graphicx}
\usepackage{tikz}
\usepackage{fancyhdr}
\usepackage{multicol}
\usepackage{datetime2}
\usepackage{hyperref}
\usepackage{mdframed}
\usepackage{amsmath}

\hypersetup{
    colorlinks=true,
    linkcolor=emergency_blue,
    citecolor=emergency_blue,
    urlcolor=emergency_blue,
    pdfborder={0 0 0}
}

% Defining page geometry
\geometry{
  top=2.5cm,
  bottom=2.5cm,
  left=2cm,
  right=2cm,
  headheight=25pt
}

% Adjusting topmargin to compensate for increased headheight
\addtolength{\topmargin}{-10pt}

% Defining colors
\definecolor{emergency_red}{RGB}{186,24,27}
\definecolor{emergency_blue}{RGB}{0,114,188}
\definecolor{emergency_gray}{RGB}{120,120,120}
\definecolor{highlight_yellow}{RGB}{255,250,205}

% Setting title formats
\titleformat{\section}[block]{\Large\bfseries\color{emergency_red}}{\thesection}{10pt}{}
\titleformat{\subsection}[block]{\large\bfseries\color{emergency_blue}}{\thesubsection}{8pt}{}
\titleformat{\subsubsection}[block]{\normalsize\bfseries\color{emergency_gray}}{\thesubsubsection}{6pt}{}

% Defining custom environment for important notes
\newmdenv[
    linecolor=emergency_red,
    backgroundcolor=highlight_yellow,
    linewidth=2pt,
    topline=true,
    bottomline=true,
    leftline=true,
    rightline=true,
    innertopmargin=10pt,
    innerbottommargin=10pt,
    innerleftmargin=10pt,
    innerrightmargin=10pt
]{importantbox}

% Defining note command with tcolorbox
\newcommand{\note}[1]{
    \par
    \noindent
    \begin{tcolorbox}[
        colback=emergency_blue!5!white,
        colframe=emergency_blue,
        title=\textbf{重點提醒},
        fonttitle=\bfseries,
        boxrule=0.5pt,
        arc=3pt,
        left=6pt,
        right=6pt,
        top=6pt,
        bottom=6pt
    ]
    #1
    \end{tcolorbox}
}

% Setting up header and footer
\pagestyle{fancy}
\fancyhf{}
\fancyhead[L]{\textcolor{emergency_red}{醫學教育文獻整理}}
\fancyhead[R]{\textcolor{emergency_blue}{謝慕揚, MD, PhD, FESC}}
\fancyfoot[C]{\textcolor{emergency_gray}{\thepage}}
\renewcommand{\headrulewidth}{1pt}
\renewcommand{\headrule}{\hbox to\headwidth{\color{emergency_red}\leaders\hrule height \headrulewidth\hfill}}

% Defining custom date format for ISO 8601
\DTMnewdatestyle{iso}{
  \renewcommand{\DTMdisplaydate}[4]{\number##1-\number##2-\number##3}
  \renewcommand{\DTMDisplaydate}[4]{\number##1-\number##2-\number##3}
}
\DTMsetdatestyle{iso}
\newcommand{\mydate}{\DTMtoday}

\begin{document}

\title{\textcolor{emergency_red}{臨床不確定性的教學策略}}
\author{\textcolor{emergency_blue}{文獻整理：謝慕揚, MD, PhD, FESC}}
\date{\textcolor{emergency_gray}{\mydate}}
\maketitle

\section{文章基本資訊}

\begin{longtable}{>{\raggedright\arraybackslash}p{3cm} >{\raggedright\arraybackslash}p{12cm}}
\toprule
\textbf{\textcolor{emergency_red}{項目}} & \textbf{\textcolor{emergency_red}{內容}} \\
\midrule
\endfirsthead

\multicolumn{2}{c}%
{{\bfseries \tablename\ \thetable{} -- 接續前頁}} \\
\toprule
\textbf{\textcolor{emergency_red}{項目}} & \textbf{\textcolor{emergency_red}{內容}} \\
\midrule
\endhead

\midrule \multicolumn{2}{r}{{接續下頁}} \\ \midrule
\endfoot

\bottomrule
\endlastfoot

期刊 & New England Journal of Medicine (NEJM) \\
文章標題 & Educational Strategies to Prepare Trainees for Clinical Uncertainty \\
作者 & Jonathan S. Ilgen, M.D., Ph.D.$^1$ 和 Gurpreet Dhaliwal, M.D.$^{2,3}$ \\
所屬機構 & $^1$Department of Emergency Medicine, University of Washington, Seattle \newline $^2$Department of Medicine, University of California, San Francisco \newline $^3$Medical Service, San Francisco Veterans Affairs Medical Center \\
發表日期 & October 23, 2025 \\
期刊卷期 & Volume 393, Number 16 \\
頁碼 & 1624-1632 \\
DOI & \href{https://doi.org/10.1056/NEJMra2408797}{10.1056/NEJMra2408797} \\
文章類型 & Review Article (Medical Education series) \\
編輯 & Jeffrey M. Drazen, M.D., Raja-Elie Abdulnour, M.D., and Judith L. Bowen, M.D., Ph.D. \\
\end{longtable}

\section{核心要點 (Key Points)}

\note{不確定性是醫療實務中無所不在的現象。如果沒有不確定性，社會就不需要醫師來做判斷。}

\begin{longtable}{>{\raggedright\arraybackslash}p{15cm}}
\toprule
\textbf{\textcolor{emergency_blue}{核心概念}} \\
\midrule
\endfirsthead

\multicolumn{1}{c}%
{{\bfseries \tablename\ \thetable{} -- 接續前頁}} \\
\toprule
\textbf{\textcolor{emergency_blue}{核心概念}} \\
\midrule
\endhead

\midrule \multicolumn{1}{r}{{接續下頁}} \\ \midrule
\endfoot

\bottomrule
\endlastfoot

管理臨床不確定性是醫師的基本技能，也是有效臨床推理的核心 \\

臨床醫師透過一系列線索來辨識不確定性，並藉由前瞻性規劃 (forward planning) 和監測 (monitoring) 的循環過程來應對這些經驗 \\

教師可以透過鼓勵討論不確定性的潛在來源（診斷、治療或預後）和類型來幫助學員管理不確定性： \newline • Epistemic uncertainty（認識論不確定性）：額外資訊可能重新框架或降低不確定性 \newline • Aleatoric uncertainty（偶然性不確定性）：結果本質上是可變的 \\

教師應該將不確定性重新定義為學習機會的訊號 \\

科技不會解決不確定性；它只是轉移了不確定性的焦點 \\
\end{longtable}

\section{不確定性與相關術語的定義}

\begin{longtable}{>{\raggedright\arraybackslash}p{4cm} >{\raggedright\arraybackslash}p{11cm}}
\toprule
\textbf{\textcolor{emergency_red}{術語}} & \textbf{\textcolor{emergency_red}{定義}} \\
\midrule
\endfirsthead

\multicolumn{2}{c}%
{{\bfseries \tablename\ \thetable{} -- 接續前頁}} \\
\toprule
\textbf{\textcolor{emergency_red}{術語}} & \textbf{\textcolor{emergency_red}{定義}} \\
\midrule
\endhead

\midrule \multicolumn{2}{r}{{接續下頁}} \\ \midrule
\endfoot

\bottomrule
\endlastfoot

Uncertainty \newline （不確定性） & 不確定或懷疑的感覺；風險（通常是無法計算的）、模糊性和複雜性都會導致這種感覺 \\

Clinical Uncertainty \newline （臨床不確定性） & 對臨床情況的理解缺乏完全的信心，常伴隨著「發生了什麼事？」或「我該怎麼做？」的想法 \\

Probability \newline （機率） & 未來事件發生的可能性 \\

Risk \newline （風險） & 傷害、損失或負面結果的可能性或機率 \\

Ambiguity \newline （模糊性） & 決策所需的資訊缺失、不精確或相互衝突，或資訊可以用兩種或更多不同的方式解釋（例如：胸部 X 光片上是否存在浸潤） \\

Complexity \newline （複雜性） & 情況中多個組成部分以多種方式交互作用，可能導致不可預測的結果，或難以理解的模式 \\
\end{longtable}

\section{臨床不確定性的經驗}

\subsection{不確定性的辨識線索}

\begin{longtable}{>{\raggedright\arraybackslash}p{6cm} >{\raggedright\arraybackslash}p{9cm}}
\toprule
\textbf{\textcolor{emergency_blue}{線索類型}} & \textbf{\textcolor{emergency_blue}{具體表現}} \\
\midrule
\endfirsthead

\multicolumn{2}{c}%
{{\bfseries \tablename\ \thetable{} -- 接續前頁}} \\
\toprule
\textbf{\textcolor{emergency_blue}{線索類型}} & \textbf{\textcolor{emergency_blue}{具體表現}} \\
\midrule
\endhead

\midrule \multicolumn{2}{r}{{接續下頁}} \\ \midrule
\endfoot

\bottomrule
\endlastfoot

\textcolor{emergency_red}{認知線索 (Cognitive Cues)} & 我不確定發生了什麼或該怎麼做 \\

\textcolor{emergency_red}{身體線索 (Somatic Cues)} & 我的胃部有「下沉」的感覺 \\

\textcolor{emergency_red}{情緒線索 (Emotional Cues)} & 這裡有什麼感覺不太對勁 \\

\textcolor{emergency_red}{行為線索 (Behavioral Cues)} & 我注意到自己放慢了速度 \\
\end{longtable}

\subsection{不確定性的兩種觀點}

\begin{longtable}{>{\raggedright\arraybackslash}p{4cm} >{\raggedright\arraybackslash}p{11cm}}
\toprule
\textbf{\textcolor{emergency_blue}{觀點}} & \textbf{\textcolor{emergency_blue}{說明}} \\
\midrule
\endfirsthead

\multicolumn{2}{c}%
{{\bfseries \tablename\ \thetable{} -- 接續前頁}} \\
\toprule
\textbf{\textcolor{emergency_blue}{觀點}} & \textbf{\textcolor{emergency_blue}{說明}} \\
\midrule
\endhead

\midrule \multicolumn{2}{r}{{接續下頁}} \\ \midrule
\endfoot

\bottomrule
\endlastfoot

\textcolor{emergency_red}{特質觀點 (Trait)} & 將醫師應對不確定性的能力視為跨越一系列經驗的穩定人格特質 \newline • 透過量表測量「不確定性耐受度」 \newline • 較低的不確定性耐受度與較高的倦怠率、較低的職業滿意度和專科選擇有關 \\

\textcolor{emergency_red}{狀態觀點 (State)} & 不確定性在當下產生，反映知識和技能與情境中臨床問題交互作用的特殊方式 \newline • 來自自我調節文獻 \newline • 描述人們如何使用來自自己和世界的線索來辨識不確定性 \newline • 對自己處理這些情況的能力做出判斷 \\
\end{longtable}

\section{經驗豐富的臨床醫師如何管理不確定性}

\begin{importantbox}
\textbf{核心概念框架}

經驗豐富的臨床醫師使用兩個關鍵過程來管理不確定性：
\begin{enumerate}
\item \textbf{前瞻性規劃 (Forward Planning)}：模擬問題可能如何演變，預測可能發生的情況，思考應變措施，並制定安全計畫
\item \textbf{監測 (Monitoring)}：注意來自患者、自己的反應和環境的更廣泛線索，持續重新審視對正在發生的事情和如何處理情況的想法
\end{enumerate}
\end{importantbox}

\subsection{前瞻性規劃 (Forward Planning)}

\begin{longtable}{>{\raggedright\arraybackslash}p{15cm}}
\toprule
\textbf{\textcolor{emergency_blue}{前瞻性規劃的關鍵問題}} \\
\midrule
\endfirsthead

\multicolumn{1}{c}%
{{\bfseries \tablename\ \thetable{} -- 接續前頁}} \\
\toprule
\textbf{\textcolor{emergency_blue}{前瞻性規劃的關鍵問題}} \\
\midrule
\endhead

\midrule \multicolumn{1}{r}{{接續下頁}} \\ \midrule
\endfoot

\bottomrule
\endlastfoot

隨著這個問題的演變，可能會發生什麼？ \\
潛在的風險是什麼？ \\
如果事情變得更糟，我應該打電話給誰？ \\
\end{longtable}

\subsection{監測 (Monitoring)}

\begin{longtable}{>{\raggedright\arraybackslash}p{4cm} >{\raggedright\arraybackslash}p{11cm}}
\toprule
\textbf{\textcolor{emergency_blue}{監測面向}} & \textbf{\textcolor{emergency_blue}{關鍵問題}} \\
\midrule
\endfirsthead

\multicolumn{2}{c}%
{{\bfseries \tablename\ \thetable{} -- 接續前頁}} \\
\toprule
\textbf{\textcolor{emergency_blue}{監測面向}} & \textbf{\textcolor{emergency_blue}{關鍵問題}} \\
\midrule
\endhead

\midrule \multicolumn{2}{r}{{接續下頁}} \\ \midrule
\endfoot

\bottomrule
\endlastfoot

\textcolor{emergency_red}{自我監測 (Self)} & 我的假設如何演變？ \newline 我對管理這種情況感到舒適嗎？ \\

\textcolor{emergency_red}{患者監測 (Patient)} & 哪些發現改變了？ \newline 患者理解治療計畫嗎？ \\

\textcolor{emergency_red}{情境監測 (Situation)} & 其他人似乎在擔心什麼？ \newline 我有合適的資源嗎？ \\
\end{longtable}

\section{教師的建議策略}

\subsection{敘述不確定性 (Narrating Uncertainty)}

\begin{longtable}{>{\raggedright\arraybackslash}p{15cm}}
\toprule
\textbf{\textcolor{emergency_red}{策略說明}} \\
\midrule
\endfirsthead

\multicolumn{1}{c}%
{{\bfseries \tablename\ \thetable{} -- 接續前頁}} \\
\toprule
\textbf{\textcolor{emergency_red}{策略說明}} \\
\midrule
\endhead

\midrule \multicolumn{1}{r}{{接續下頁}} \\ \midrule
\endfoot

\bottomrule
\endlastfoot

\textbf{核心概念：}透過大聲思考臨床問題，教師可以敘述不確定性帶來挑戰的方式 \\

\textbf{具體做法：}\newline • 分享「我不確定診斷」或「這裡有些事情感覺不對勁」的想法 \newline • 描述如何利用這些反應進行規劃和監測 \newline • 這種例行性的披露可以將不確定性正常化 \\

\textbf{不確定性檢查 (Uncertainty Check)：}\newline 詢問學員對診斷或計畫的信心程度，是一種簡單的方法來揭示臨床不確定性並促進討論 \\

\textbf{「如果...那麼」計畫 (If-Then Plans)：}\newline 有效方法來強調學員應該最關注案例的哪些方面，以及應該在哪些停止點尋求協助 \newline 例如：「雖然靜脈注射抗生素治療腎盂腎炎的前1-2天發燒是典型的，但我對72小時時仍發燒較不放心——那時我會安排CT掃描以排除腎結石或腎膿瘍」 \\
\end{longtable}

\subsection{框架不確定性 (Framing Uncertainty)}

\begin{longtable}{>{\raggedright\arraybackslash}p{4cm} >{\raggedright\arraybackslash}p{11cm}}
\toprule
\textbf{\textcolor{emergency_blue}{框架類型}} & \textbf{\textcolor{emergency_blue}{說明與應用}} \\
\midrule
\endfirsthead

\multicolumn{2}{c}%
{{\bfseries \tablename\ \thetable{} -- 接續前頁}} \\
\toprule
\textbf{\textcolor{emergency_blue}{框架類型}} & \textbf{\textcolor{emergency_blue}{說明與應用}} \\
\midrule
\endhead

\midrule \multicolumn{2}{r}{{接續下頁}} \\ \midrule
\endfoot

\bottomrule
\endlastfoot

\textcolor{emergency_red}{位置框架} & 不確定性涉及診斷、治療、預後，還是工作場所動態？ \\

\textcolor{emergency_red}{Epistemic Uncertainty \newline （認識論不確定性）} & 有限知識或資訊的問題，可以被解決 \newline 例如：「是否有證據支援在攝入cocaine的患者中使用beta-blocker的安全性？」 \newline 應對：提供資訊或分享過去經驗 \\

\textcolor{emergency_red}{Aleatoric Uncertainty \newline （偶然性不確定性）} & 反映生物醫學和健康系統中不可簡化的隨機性 \newline 例如：「患者可能會有藥物不良反應嗎？」 \newline 應對：評估事件發生的機率或思考較不可預測或最壞情況的應變措施 \\
\end{longtable}

\subsection{模擬不確定性 (Simulating Uncertainty)}

\begin{longtable}{>{\raggedright\arraybackslash}p{15cm}}
\toprule
\textbf{\textcolor{emergency_red}{策略說明}} \\
\midrule
\endfirsthead

\multicolumn{1}{c}%
{{\bfseries \tablename\ \thetable{} -- 接續前頁}} \\
\toprule
\textbf{\textcolor{emergency_red}{策略說明}} \\
\midrule
\endhead

\midrule \multicolumn{1}{r}{{接續下頁}} \\ \midrule
\endfoot

\bottomrule
\endlastfoot

\textbf{傳統方法的限制：}病例討論和模擬通常被設計成線性的、邏輯的和可解決的，以傳達診斷和管理的基礎知識 \\

\textbf{進階訓練的轉變：}在更高級的訓練階段，教師可以將病例討論轉向具有一系列合理選項和多個無法解決的困境的情境 \\

\textbf{不確定性情境案例：}例如，急性單關節炎的病例，發現與晶體性和感染性關節炎都相符，會促使學習者在沒有明確診斷的情況下應對管理的各個方面 \newline 討論將針對學員如何實時定義和應對這些困境，例如： \newline • 滑膜液中細胞計數閾值的非確定性 \newline • 在深夜聯繫顧問的不安 \newline • 對患者是否負擔得起處方治療的擔憂 \\

\textbf{高擬真度模擬：}使用高擬真度人體模型或互動式線上病例，重現臨床醫師在實務中面臨的各種不確定性 \newline 教師可以邀請學員「凍結時間」並討論： \newline • 他們在自己和他人身上注意到的線索 \newline • 描述他們對情況演變的理解 \newline • 討論一系列管理方法 \\

\textbf{教師角色轉變：}不是提供已知問題的答案（例如「這個患者有主動脈剝離，其管理包括...」），教師敘述他們的方法作為教學工具（例如「我曾經有過這樣棘手的病例，以下是當我努力理清楚發生了什麼時幫助我的幾種策略」） \\
\end{longtable}

\subsection{溝通不確定性 (Communicating Uncertainty)}

\begin{longtable}{>{\raggedright\arraybackslash}p{4cm} >{\raggedright\arraybackslash}p{11cm}}
\toprule
\textbf{\textcolor{emergency_blue}{溝通對象}} & \textbf{\textcolor{emergency_blue}{工具與方法}} \\
\midrule
\endfirsthead

\multicolumn{2}{c}%
{{\bfseries \tablename\ \thetable{} -- 接續前頁}} \\
\toprule
\textbf{\textcolor{emergency_blue}{溝通對象}} & \textbf{\textcolor{emergency_blue}{工具與方法}} \\
\midrule
\endhead

\midrule \multicolumn{2}{r}{{接續下頁}} \\ \midrule
\endfoot

\bottomrule
\endlastfoot

\textcolor{emergency_red}{向督導者溝通} & 病例呈報模型嵌入學生提問的步驟（例如「我想知道為什麼患者仍然咳嗽」），讓他們有機會詢問督導者關於不確定性的問題 \\

\textcolor{emergency_red}{向同事溝通} & 在照護交接期間使用交班工具可以促使臨床醫師與同事分享診斷不確定性的程度（一些不確定性、顯著不確定性或高度不確定性） \\

\textcolor{emergency_red}{向患者溝通} & 多種資源可用於教導學員如何向患者溝通不確定性，包括學習如何描述它以及如何評估患者對此類訊息的接受度 \newline 研究顯示：家屬偏好包含鑑別診斷的診斷不確定性溝通（「疼痛最可能是踝關節扭傷引起的，儘管痛風是一種可能性」），而不是更一般的不確定性表達（「我不確定是什麼引起的疼痛」） \\
\end{longtable}

\section{臨床教學的預診策略總表}

\begin{longtable}{>{\raggedright\arraybackslash}p{3cm} >{\raggedright\arraybackslash}p{4cm} >{\raggedright\arraybackslash}p{4cm} >{\raggedright\arraybackslash}p{4cm}}
\toprule
\textbf{\textcolor{emergency_red}{技能}} & \textbf{\textcolor{emergency_red}{情境或線索}} & \textbf{\textcolor{emergency_red}{教育診斷}} & \textbf{\textcolor{emergency_red}{教育策略}} \\
\midrule
\endfirsthead

\multicolumn{4}{c}%
{{\bfseries \tablename\ \thetable{} -- 接續前頁}} \\
\toprule
\textbf{\textcolor{emergency_red}{技能}} & \textbf{\textcolor{emergency_red}{情境或線索}} & \textbf{\textcolor{emergency_red}{教育診斷}} & \textbf{\textcolor{emergency_red}{教育策略}} \\
\midrule
\endhead

\midrule \multicolumn{4}{r}{{接續下頁}} \\ \midrule
\endfoot

\bottomrule
\endlastfoot

\textcolor{emergency_blue}{敘述不確定性} & 學員呈報的評估看似完美且明確，沒有提及來自病史和體檢的不確定性 & 學員可能為了表現出更大的確定性和信心而過度簡化 & 重新整合不一致的病史或體檢發現以提出工作假設的疑問 \newline 例：「我想知道這個皮疹如何改變我們的思考。對我們的工作診斷有不符合的事情是正常的」 \\

\textcolor{emergency_blue}{框架不確定性} & 學員呈報廣泛的不確定性，從診斷和管理的具體問題到關於證據或實務變異的更廣泛問題 & 學員可能受益於組織這些不確定性的策略 & 幫助學員將問題分層為認識論（額外資訊可能重新框架或降低不確定性）或偶然性（結果本質上是可變的） \newline 例：「您強調了關於這個病例的幾個重要問題。我們應該在哪裡尋求資訊以降低此時的不確定性？這個病例的哪些方面無論我們做什麼都可能仍然未知或不可預測？」 \\

\textcolor{emergency_blue}{前瞻性規劃} & 學員結束關於患者的討論而沒有規劃不良結果的應變措施 & 學員可能沒有教師心中的相同鑑別診斷，或可能不分享對風險的相同關切，也可能缺乏該問題的經驗 & 挑戰學員思考病例可能展開的不同方式以及學員將如何應對 \newline 例：「我們正在制定的計畫可能有哪三種不同的潛在結果？您現在可以做些什麼來為這些可能性做準備？」 \\

\textcolor{emergency_blue}{患者監測} & 學員結束關於患者的討論而沒有具體的監測計畫 & 學員可能不知道隨著病例演變或開始治療應該尋找哪些發現 & 強調要監測的關鍵發現，以確定診斷和管理是否在正軌上 \newline 例：「您計畫何時再次評估這位患者？您在尋找什麼樣的訊號？發燒的出現如何改變您的思考或管理？」 \\

\textcolor{emergency_blue}{自我監測} & 學員在病例呈報期間表達擔憂，或教師注意到表明學員可能對病例感到擔心或焦慮的行為 & 學員可能沒有將情緒或身體線索視為有用的不確定性訊號 & 探究學員的反應並驗證它們作為可用於管理不確定性的有價值訊號 \newline 例：「我在您的病例呈報中聽到您說『擔心』，我對這個病例也有同樣的反應。讓我們比較一下讓我們每個人不確定的因素以及它們為什麼引起關切」 \\

\textcolor{emergency_blue}{監測他人} & 學員忽視其他醫療專業人員的行為或語言，這些行為或語言表明他們對情況的思考不同 & 學員可能不知道特定醫療團隊如何共同工作，包括其他人試圖表達關切的時刻 & 分享您對其他醫療專業人員使用的行為和語言的解釋 \newline 例：「護理師的語氣和頻繁更新表明對這位患者的擔憂。您可以使用什麼策略來探索護理師的思考？」 \\

\textcolor{emergency_blue}{溝通不確定性} & 教師注意到學員沒有向患者披露不確定性 & 學員可能擔心揭示不確定性可能會破壞患者對學員技能的信心 & 示範如何與患者分享不確定性，包括邀請患者幫助指導監測計畫 \newline 例：「注意我如何通過說明我們認為診斷是痛風但我們將繼續考慮蜂窩組織炎作為替代診斷來開始對話。然後我會詢問患者分享他的理解和關切。患者和我將討論一系列選項，以幫助建立一個我們都可以接受的計畫」 \\
\end{longtable}

\section{重要引文與概念}

\note{「如果沒有不確定性，社會就不需要醫師來做判斷。」}

\begin{longtable}{>{\raggedright\arraybackslash}p{15cm}}
\toprule
\textbf{\textcolor{emergency_blue}{關鍵概念與語句}} \\
\midrule
\endfirsthead

\multicolumn{1}{c}%
{{\bfseries \tablename\ \thetable{} -- 接續前頁}} \\
\toprule
\textbf{\textcolor{emergency_blue}{關鍵概念與語句}} \\
\midrule
\endhead

\midrule \multicolumn{1}{r}{{接續下頁}} \\ \midrule
\endfoot

\bottomrule
\endlastfoot

「不確定性在醫療實務中無所不在，但通常被視為需要容忍或根除的令人遺憾的現象。為了讓學員做好實務準備，教師需要將不確定性重新定義為理解和管理臨床問題的核心特徵。」 \\

「對不確定性的厭惡會導致潛意識地努力在日常臨床工作中隱藏、抑制或最小化它。然而，即使隨著技術和治療的進步，不確定性的邊界發生了變化，不確定性也無法從臨床實務中根除。」 \\

「有效的教學傳達出不確定性不應該被避免，而是實務的現實，可以被識別和管理。」 \\

「學員對臨床不確定性的看法與他們自己的看法不同，即使在照護同一位患者時也是如此。這些差異源於學員正在發展的知識或技能、對周圍設備或跨專業團隊結構的不熟悉，或在實務中與類似情況的有限經驗。」 \\

「臨床教師在指出和命名學員的不確定性方面發揮著關鍵作用，並向他們展示即使對臨床情況的理解不完整，如何安全地前進。」 \\

「教師應該將不確定性重新定義為有學習機會的訊號。」 \\

「科技不會解決不確定性；它只是轉移了不確定性的焦點。」 \\
\end{longtable}

\section{學員的典型不確定性案例}

\begin{importantbox}
\textbf{呼吸困難患者的假設案例}

一位第一年住院醫師在門診評估一位呼吸困難的患者時，可能會有一系列的不確定性：

\begin{itemize}
\item 他是否真的檢測到頸靜脈怒張？
\item 應該選擇哪些初始治療？
\item 如果患者的臨床狀況突然惡化，需要什麼資源？
\item 在哪裡可以找到幫助？
\item 他的不適是遇到陌生問題時訓練的預期部分，還是這些反應表明他的督導者也會分享的關切？
\end{itemize}
\end{importantbox}

\section{參考文獻選錄}

\begin{longtable}{>{\raggedright\arraybackslash}p{2cm} >{\raggedright\arraybackslash}p{13cm}}
\toprule
\textbf{\textcolor{emergency_red}{編號}} & \textbf{\textcolor{emergency_red}{文獻}} \\
\midrule
\endfirsthead

\multicolumn{2}{c}%
{{\bfseries \tablename\ \thetable{} -- 接續前頁}} \\
\toprule
\textbf{\textcolor{emergency_red}{編號}} & \textbf{\textcolor{emergency_red}{文獻}} \\
\midrule
\endhead

\midrule \multicolumn{2}{r}{{接續下頁}} \\ \midrule
\endfoot

\bottomrule
\endlastfoot

1 & Simpkin AL, Schwartzstein RM. Tolerating uncertainty — the next medical revolution? N Engl J Med 2016;375:1713-5. \\

2 & Djulbegovic B, Hozo I, Greenland S. Uncertainty in clinical medicine. In: Gifford F, ed. Philosophy of Medicine. Amsterdam: North Holland, 2011:299-356. \\

3 & Ilgen JS, Eva KW, Regehr G. What's in a label? Is diagnosis the start or the end of clinical reasoning? J Gen Intern Med 2016;31:435-7. \\

11 & Ilgen JS, Eva KW, de Bruin A, Cook DA, Regehr G. Comfort with uncertainty: reframing our conceptions of how clinicians navigate complex clinical situations. Adv Health Sci Educ Theory Pract 2019;24:797-809. \\

19 & Ilgen JS, Bowen JL, de Bruin ABH, Regehr G, Teunissen PW. "I Was Worried About the Patient, but I Wasn't Feeling Worried": how physicians judge their comfort in settings of uncertainty. Acad Med 2020;95:Suppl:S67-S72. \\

21 & Ilgen JS, Teunissen PW, de Bruin ABH, Bowen JL, Regehr G. Warning bells: how clinicians leverage their discomfort to manage moments of uncertainty. Med Educ 2021;55:233-41. \\

22 & Moulton C, Regehr G, Lingard L, Merritt C, Macrae H. 'Slowing down when you should': initiators and influences of the transition from the routine to the effortful. J Gastrointest Surg 2010;14:1019-26. \\

24 & Eva KW. What every teacher needs to know about clinical reasoning. Med Educ 2005;39:98-106. \\

44 & Wray CM, Loo LK. The diagnosis, prognosis, and treatment of medical uncertainty. J Grad Med Educ 2015;7:523-7. \\

51 & Stephens GC, Lazarus MD. Twelve tips for developing healthcare learners' uncertainty tolerance. Med Teach 2024;46:1035-43. \\
\end{longtable}

\section{文獻整理者總結}

\begin{importantbox}
\textbf{文獻整理：謝慕揚, MD, PhD, FESC}

本文是 NEJM 醫學教育系列的優秀回顧文章，深入探討了臨床不確定性的本質及其教學策略。文章最大的價值在於將不確定性從「需要被消除的缺陷」重新定義為「臨床推理的核心特徵」，並提供了具體可行的教學方法。

\textbf{核心洞見：}

\begin{enumerate}
\item \textbf{典範轉移：}將不確定性從問題轉變為機會，這種思維方式的轉變對於醫學教育至關重要
\item \textbf{實務框架：}前瞻性規劃與監測的雙循環模型為理解和教導臨床決策提供了清晰的架構
\item \textbf{教學工具：}四大策略（敘述、框架、模擬、溝通）提供了完整的教學工具箱
\item \textbf{情境依賴：}認識到學員的不確定性經驗受其發展階段和情境脈絡影響，不應簡單視為能力不足
\end{enumerate}

\textbf{對臨床教師的啟示：}

本文提供的教學策略表格特別實用，可直接應用於床邊教學。教師應該：
\begin{itemize}
\item 主動分享自己的不確定性經驗，打破「專家永遠確定」的迷思
\item 使用「不確定性檢查」和「如果...那麼」計畫作為常規教學工具
\item 區分認識論不確定性（可透過學習解決）與偶然性不確定性（需要應變規劃）
\item 在模擬和病例討論中刻意設計包含不確定性的情境
\end{itemize}

\textbf{對心臟科教學的啟示：}

在心血管急症處理中，不確定性管理尤其重要。例如：
\begin{itemize}
\item 急性胸痛的鑑別診斷常需在不完整資訊下做出關鍵決策
\item 心導管併發症的處理需要即時的前瞻性規劃和持續監測
\item 複雜心衰竭患者的治療調整充滿偶然性不確定性
\end{itemize}

本文提供的框架可幫助我們更有系統地教導住院醫師在這些高風險情境中管理不確定性。

整理日期：\mydate, with assistance by Claude.ai
\end{importantbox}

\end{document}